\section{Discusión}
Los resultados del sistema de gestión de tickets se alinean estrechamente con lo planteado en la literatura. Laravel, como framework backend, demostró ser no solo accesible, sino también seguro y escalable, tal como destacan Subecz (2021) y Alrahhal & Alzoubi (2021). Aunque CodeIgniter mostró rendimientos similares en estudios comparativos (Universidad Muhammadiyah Surakarta, 2021), Laravel sobresale por su ecosistema de paquetes y su comunidad activa, aspectos clave para garantizar soporte y actualizaciones continuas.

En cuanto al frontend, React se consolidó como la mejor opción para construir una interfaz moderna y eficiente, respaldado por investigaciones que lo destacan frente a Angular y Vue en términos de flexibilidad (Kaur & Kaur, 2020; Singh & Kaur, 2020). La posibilidad de evolucionar hacia React Native fortalece la proyección multiplataforma, coincidiendo con Silva et al. (2022).

Desde la perspectiva metodológica, la adopción de Scrum resultó fundamental para gestionar cambios en los requerimientos y mantener un ritmo de entrega constante. Tal como lo plantean Alarcón et al. (2021), la efectividad de Scrum depende de combinarlo con buenas prácticas de documentación y pruebas. Esto se reflejó en el proyecto, donde la integración de testing automatizado contribuyó a la reducción de defectos.

Finalmente, los aportes de artículos sobre automatización, accesibilidad y seguridad (Redalyc, 2020; Redalyc, 2019) permitieron enriquecer el enfoque del proyecto, confirmando que el desarrollo de software moderno debe trascender la simple implementación técnica y contemplar la inclusión, la calidad y la sostenibilidad.